\begin{description}
\item[(T09.1)]
Angular radius of the diffraction disk is given by
$\dfrac{1.22\lambda}{D}$. The radius of the image is therefore,
\[ d_{\mathrm{image}} = f\left(\frac{1.22\lambda}{D}\right) \]
% sic: the source says "The radius of the image is therefore" although
% d_image is the diameter asked for and the boxed value is the diameter.
\hfill(0.5 marks)
\[ d_{\mathrm{image}} = 8.95 \times 10^{-6}\,\mathrm{m} \]
\hfill(0.5 marks)

\item[(T09.2a)]
% FIGURE NEEDED: hand-drawn ray diagram at the tilted boundary (tilt theta
% to the horizontal) between media n_2 (above) and n_1 (below): vertical
% incident ray, normal to the boundary, angle of incidence theta, refracted
% ray making angle (theta - alpha) with the normal and angular shift alpha
% from the vertical; 2025_theory_solution.pdf page 31.
The ray-diagram at the boundary shows the incident and refracted rays with
the normal to the boundary: the angle of incidence is $\theta$, the
refracted ray in $n_1$ makes angle $(\theta - \alpha)$ with the normal, and
$\alpha$ is the angular shift of the refracted ray from the incident
direction. \hfill(2.0 marks)

\item[(T09.2b)]
Using Snell's Law,
\begin{gather*}
n_2\sin\theta = n_1\sin(\theta - \alpha) \\
\frac{n_2}{n_1}\sin\theta = \sin\theta\cos\alpha - \cos\theta\sin\alpha
\end{gather*}
\hfill(1.0 marks)\\
Using small angle approximation ($\sin\alpha \approx \alpha$ and
$\cos\alpha \approx 1$)
\[ \frac{n_2}{n_1}\sin\theta = \sin\theta - \alpha\cos\theta \]
\[ \alpha = \frac{\sin\theta}{\cos\theta}\left(1 - \frac{n_2}{n_1}\right) \]
\hfill(1.0 marks)

\item[(T09.2c)]
The shift in position of image due to change in $\theta$ is:
\begin{align*}
\Delta x_\theta &= f\alpha
  = f\Delta\tan\theta\left(1 - \frac{n_2}{n_1}\right)
  && \text{(1.0 marks)} \\
  &= f\left(1 - \frac{n_2}{n_1}\right)
     (\tan(\theta + \Delta\theta) - \tan\theta) \\
  &= f\left(1 - \frac{n_2}{n_1}\right)\sec^2\theta\,\Delta\theta
  && \text{(1.0 marks)} \\
  &= f \times (6.98 \times 10^{-8}\,\mathrm{rad})
\end{align*}
\[ \Delta x_\theta = 1.40 \times 10^{-7}\,\mathrm{m} \]
\hfill(1.0 marks)

\item[(T09.2d)]
Shift in the position of image due to the change in $n_2$ is:
\begin{align*}
\Delta x_n &= f\tan\theta
  \left(-\frac{n_2 + \Delta n_2}{n_1} + \frac{n_2}{n_1}\right) \\
  &= -f\tan\theta\left(\frac{0.0001}{100} \times \frac{n_2}{n_1}\right)
  && \text{(2.0 marks)} \\
  &= -f \times (5.77 \times 10^{-7}\,\mathrm{rad})
\end{align*}
\[ \Delta x_n = -1.15 \times 10^{-6}\,\mathrm{m} \]
\hfill(1.0 marks)

\item[(T09.3)]
Blue light has larger refractive index than the red light, hence it will
bend more. So incoming light will vary in colour from red to blue towards
right.

So, the correct answer will be F. \hfill(2.0 marks)

Note: convex lens inverts image, but does not affect refractive index
dependent bending.

\item[(T09.4a)]
% FIGURE NEEDED: hand-drawn continuation of the zigzag-boundary diagram --
% three sets of parallel refracted rays below the boundary, with the rays
% from adjacent upward-bending planes diverging so that a patch marked "X"
% on the telescope plane (gray dotted line) receives no light;
% 2025_theory_solution.pdf page 34.
The paths of the incident light rays are drawn below the zigzag boundary:
each planar segment produces a set of parallel refracted rays, and just
below each point where two planes make an upward bend the rays diverge,
leaving a region marked "X" on the plane of the telescope objective where no
light rays will reach. \hfill(4.0 marks)

\item[(T09.4b)]
Each set of planes, with specific tilt, will lead to one set of stripes of
parallel rays. From solution of part (T09.2b),
\[ \alpha = \tan\theta\,\frac{(n_1 - n_2)}{n_1} \]
Two sets of planes with relative angle of $2\theta$ will create angular
separation of $2\alpha$ between the stripes. \hfill(1.0 marks)\\
At a distance of $H$ this will create a horizontal separation of
$\approx 2\alpha H$. For $H = 1$\,km, $\theta = 10^\circ$, the horizontal
shift will be
\[ \Delta x \simeq 2\alpha H = 2\tan\theta\,\frac{(n_1 - n_2)}{n_1}H
   = 2 \times 1.76 \times 10^{-6} \times 1 \times 10^{5}\,\mathrm{cm}
   = 0.35\,\mathrm{cm}. \]
\hfill(1.0 marks)\\
Just below the point where two planes make upward bend, a patch of this
width 0.35\,cm will be centred where no light from the star will reach.
Hence, $W_{\mathrm{X}} = 0.35$\,cm. \hfill(1.0 marks)

\item[(T09.4c)]
From the figure one can see that two adjacent planes (with downward bend),
will lead to overlapping stripes of rays producing two images.

Therefore, the region in which only single image is formed will be 10.0\,cm
(including the region X). Hence, $D_{\mathrm{max}} = 10.0$\,cm.
\hfill(4.0 marks)

\item[(T09.5)]
For $D \gg d$, several stripes will fall on the lens, leading to multiple
images. Here, four images will form, two each from the set of planes with
zigzag shapes along $x$-axis and $y$-axis respectively.

First consider the zigzag shape of the boundary along the $x$ axis. From
solution of (T09.4), angle between these two sets of rays will be
$3.52 \times 10^{-6}$\,rad. With $f = 2$\,m, the separation between the two
images will be $= 7.0\,\mu\mathrm{m}$.

For $D = 100$\,cm, the diffraction peak size $= 1.34\,\mu\mathrm{m}$. So
the images formed by the two sets of planes (with tilts of $\pm 10^\circ$
from horizontal for each set) will form two spots, clearly separated along
the $x$-axis. Similarly for the $y$-axis.

The resulting speckles will be as shown below. \hfill(6.0 marks)
% FIGURE NEEDED: hand-drawn speckle pattern -- four circular,
% non-overlapping spots placed symmetrically on the x and y axes (one on
% each side of the origin); 2025_theory_solution.pdf page 35.

\item[(T09.6)]
A single image is formed when only one stripe of rays is allowed to fall on
the objective. This is only possible when the objective is placed suitably
below the zigzag pattern.

The diffraction-limited image size for a $D = 8$\,cm lens is
$8.4 \times 10^{-6}$\,m. \hfill(1.0 marks)

The angle of the atmospheric layer changes from $+10^\circ$ to $-10^\circ$\\
$\implies$ Change in angle of rays is $3.52 \times 10^{-6}$\,rad\\
$\implies$ Shift in image position $= 3.52\,\mu\mathrm{m}$.
\hfill(1.0 marks)

Now, the shift in the image position allowed is 1\,\% of the image size.\\
$\therefore$ Allowed shift $= 0.084\,\mu\mathrm{m}$ \hfill(1.0 marks)

Longest exposure time $=$ Time taken by the image to shift by
$0.084\,\mu\mathrm{m}$

Image shifts by $3.52\,\mu\mathrm{m}$ in 1\,s $\implies$
$0.084\,\mu\mathrm{m}$ shift occurs in 0.024\,s.
\[ t_{\mathrm{max}} = 0.024\,\mathrm{s} \]
\hfill(2.0 marks)
\end{description}
